% Auto-generated by Exercise5.jl — do not edit by hand.
\begin{table}[htbp]
  \centering
  \caption{Composition and Cumulative Aggregate Labor Productivity Growth}
  \label{tab:ex5_decomposition}
  \begin{threeparttable}
  \centering
  \begin{tabular}{l S[table-format=-3.1] S[table-format=-3.1] S[table-format=-2.1]}
    \toprule
     & {Baseline} & {Full shutdown} & {Difference} \\
    \midrule
    \multicolumn{4}{l}{\textit{Panel A. Cumulative aggregate productivity growth}} \\
    I.\ Within-sector productivity & 160.4 & 172.6 & 12.2 \\
    II.\ Labor reallocation & -2.1 & -2.5 & -0.4 \\
    III.\ Value-added conversion & 0.0 & 0.1 & 0.1 \\
    \midrule
    Total & 158.3 & 170.2 & 11.9 \\
    \addlinespace
    \multicolumn{4}{l}{\textit{Panel B. Decomposition of the total difference}} \\
    I.\ Accounting component & {--} & {--} & 6.9 \\
    II.\ Productivity feedback & {--} & {--} & 5.0 \\
    \midrule
    Total & {--} & {--} & 11.9 \\
    \addlinespace
    \multicolumn{4}{l}{\textit{Panel C. Decomposition of productivity feedback}} \\
    \hspace{1em}(a) Endogenous innovation & {--} & {--} & 3.3 \\
    \hspace{1em}(b) Relative-price effect & {--} & {--} & 1.3 \\
    \hspace{1em}(c) Machine market power & {--} & {--} & 0.4 \\
    \midrule
    Total & {--} & {--} & 5.0 \\
    \bottomrule
  \end{tabular}
  \begin{tablenotes}[flushleft]
    \footnotesize
    \RaggedRight
    \item[] \textit{Notes:} Entries are cumulative log-point contributions over 1950--2023.
  Difference is full shutdown minus baseline. Panel A follows Equation (23). Panel B separates
  changes in sectoral productivity growth from reweighting using Equation (24); the accounting
  component combines reweighting with the change in the labor-reallocation term. Panel C applies
  the sectoral productivity decomposition in Equation (18) to the productivity-feedback component;
  embodied technical progress and labor efficiency follow the same exogenous paths in both
  equilibria and contribute nothing. Entries are rounded so that they add to the printed totals.
  Differences and totals are calculated before rounding.
  \end{tablenotes}
  \end{threeparttable}
\end{table}
